%% Intermediate theorem logics and varieties of Heyting algebras.

\subsection{Intermediate logics and Heyting varieties}
\label{sec:p10-intermediate-algebras}

Work in the fixed algebraic signature
\(\tau_H=(\wedge,\vee,\Rightarrow,0,1)\).
The residual \(\Rightarrow\) is the operation of
Definition~\ref{def:p10-heyting-algebra}.
Write \(\mathsf{HA}\) for all Heyting algebras, including the trivial
one, and \(\mathsf{BA}\) for those whose bounded lattice reduct is
Boolean, with \(a\Rightarrow b=\neg a\vee b\).
Thus both classes are considered in the same signature.

\begin{hbtheorem}{thm:p10-heyting-equational}{proved=annex:p10.intermediate.1}{The equational class of Heyting algebras}
A bounded distributive lattice equipped with a binary operation
\(\Rightarrow\) is a Heyting algebra exactly when it satisfies
\[
\begin{aligned}
 a\Rightarrow a&=1,\\
 a\wedge(a\Rightarrow b)&=a\wedge b,\\
 b\wedge(a\Rightarrow b)&=b,\\
 a\Rightarrow(b\wedge c)&=(a\Rightarrow b)\wedge(a\Rightarrow c).
\end{aligned}
\]
Together with the bounded distributive lattice identities, these
are a finite equational axiomatization of \(\mathsf{HA}\).
Inside \(\mathsf{HA}\), the additional identity
\[
 a\vee(a\Rightarrow0)=1
\]
defines precisely \(\mathsf{BA}\).
\end{hbtheorem}

The finite residual axioms are
\cite[Theorem~2.2.6, pp.~21--22]{Bezhanishvili2006Thesis};
the Boolean characterization is Proposition~2.2.9 on p.~23.
The proof here includes both directions, so the use of varieties does
not depend on treating residuation as an equational axiom without
justification.

\begin{hbdefinition}{def:p10-intermediate-logic}{Intermediate theorem logics and their quotients}
Fix a countably infinite variable set \(X=\{p_0,p_1,\ldots\}\).
Use formulas built from \(X\) by \(\bot,\wedge,\vee,\to\), with
\(\neg\varphi=(\varphi\to\bot)\), \(\top=(\bot\to\bot)\), and
\(\varphi\leftrightarrow\psi=(\varphi\to\psi)\wedge(\psi\to\varphi)\).
IPC is the theorem set of Definition~\ref{def:p10-ipc-lindenbaum},
restricted to \(X\), and CPL is the classical theorem set in this
same grammar.
An \emph{intermediate logic} is a set of formulas \(L\) such that
\[
 \mathsf{IPC}\subseteq L\subseteq\mathsf{CPL},
\]
closed under modus ponens and uniform substitution.
Write \(\vdash_L\varphi\) for \(\varphi\in L\).
These are closure rules for theorem sets.
If local premises are used, a derivation from \(\Gamma\) is a finite
sequence using \(\Gamma\), theorems of \(L\), and modus ponens;
substitution does not act freely on premise-dependent conclusions.

Put
\[
 \varphi\equiv_L\psi\quad\Longleftrightarrow\quad
                    \vdash_L\varphi\leftrightarrow\psi,
 \qquad H_L=\operatorname{Fm}(X)/{\equiv_L}.
\]
The proposed operations on \(H_L\) are induced by the formula
connectives, with algebraic \(\Rightarrow\) induced by \(\to\), and
the proposed order is
\[
 [\varphi]_L\leq[\psi]_L
       \quad\Longleftrightarrow\quad\vdash_L\varphi\to\psi.
\]
The following theorem verifies these constructions.
\index{intermediate logic}\index{Lindenbaum--Tarski algebra!intermediate logic}
\end{hbdefinition}

\begin{hbdefinition}{def:p10-intermediate-variety-maps}{The logic--variety assignments}
For an intermediate logic \(L\), define
\[
 \mathsf V(L)=\{H\in\mathsf{HA}:H\models\lambda
                                      \text{ for every }\lambda\in L\},
\]
where validity means value \(1\) under every valuation \(X\to H\).
For a class \(\mathcal V\subseteq\mathsf{HA}\), put
\[
 \operatorname{Log}(\mathcal V)
    =\{\varphi:H\models\varphi\text{ for every }H\in\mathcal V\}.
\]
A subvariety of \(\mathsf{HA}\) is the class of Heyting algebras
satisfying some set of identities in \(\tau_H\).
Identities may be taken over the fixed infinite variable set \(X\):
every identity uses only finitely many variables, which can be renamed
to \(p_0,p_1,\ldots\) in order of first occurrence.
\index{variety!Heyting algebras}\index{algebraic semantics!intermediate logic}
\end{hbdefinition}

\begin{hbtheorem}{thm:p10-intermediate-algebra-completeness}{proved=annex:p10.intermediate.2}{Algebraic completeness of an intermediate logic}
For every intermediate \(L\), \(\equiv_L\) is a congruence,
\(H_L\) is a nontrivial Heyting algebra, and \(H_L\in\mathsf V(L)\).
For every formula \(\varphi\),
\[
 \varphi\in L
 \quad\Longleftrightarrow\quad [\varphi]_L=1
 \quad\Longleftrightarrow\quad H_L\models\varphi
 \quad\Longleftrightarrow\quad\mathsf V(L)\models\varphi.
\]
The valuation \(p\mapsto[p]_L\) refutes every non-theorem.
\end{hbtheorem}

\begin{hbtheorem}{thm:p10-intermediate-variety-correspondence}{proved=annex:p10.intermediate.3}{The lattice correspondence for intermediate logics}
The assignments \(L\mapsto\mathsf V(L)\) and
\(\mathcal V\mapsto\operatorname{Log}(\mathcal V)\) are mutually
inverse order-reversing bijections between intermediate logics and
the interval of subvarieties
\[
 \mathsf{BA}\subseteq\mathcal V\subseteq\mathsf{HA}.
\]
In particular,
\[
 \operatorname{Log}(\mathsf V(L))=L,\qquad
 \mathsf V(\operatorname{Log}(\mathcal V))=\mathcal V.
\]
The endpoints correspond as
\[
 \mathsf V(\mathsf{IPC})=\mathsf{HA},\qquad
 \mathsf V(\mathsf{CPL})=\mathsf{BA},
 \qquad
 \operatorname{Log}(\mathsf{HA})=\mathsf{IPC},\quad
 \operatorname{Log}(\mathsf{BA})=\mathsf{CPL}.
\]
These are complete lattices.  On the logic side, nonempty meets are
intersections and joins are the least MP- and substitution-closed
extensions of IPC containing the union.  The empty meet is CPL and
the empty join is IPC.
On the variety side, nonempty meets are intersections and joins are
the least subvarieties containing \(\mathsf{BA}\) and the union.
The empty meet is \(\mathsf{HA}\) and the empty join is \(\mathsf{BA}\).
The correspondence interchanges all these meets and joins.
\end{hbtheorem}

The source defines intermediate logics in Definition~2.1.4, p.~12,
states their algebraic completeness in Theorem~2.2.18, p.~25, and
gives the lattice anti-isomorphism in Theorem~2.2.19, p.~26
\cite{Bezhanishvili2006Thesis}.
The interval above makes the consistent-logic convention explicit:
the trivial algebra belongs to every variety in it, but the variety
consisting only of trivial algebras is outside this interval.
These are algebraic completeness statements; no assertion of ordinary
Kripke completeness for every intermediate logic follows from them.
